% COMBINED Letter of Intent (LoI) + appendix - Marta Parazzini (CNR-IEIIT)
% to Robin Wydaeghe & Prof. Wout Joseph. One signable document:
%   page 1 = the LoI (explicitly headed "Letter of Intent (LoI)"), which references
%   the Appendix = the 22/06/2026 record of interview with her verbatim quotes.
% Title corrected to "Director of Research". Licence phrasing generalised (not "a licence").
% Disclosure guard: method only "substantially faster than FDTD, similar accuracy, differentiable".
% Compile: pdflatex parazzini_letter_of_intent_combined.tex
\documentclass[11pt]{article}
\usepackage[a4paper,margin=2.5cm]{geometry}
\usepackage{parskip}
\usepackage{graphicx}
\usepackage{enumitem}
\usepackage[T1]{fontenc}
\usepackage{lmodern}
\pagestyle{empty}

\begin{document}

% --- CNR-IEIIT letterhead placeholder ---------------------------------------
% \includegraphics[height=1.6cm]{cnr_ieiit_logo}\\[0.5em]
{\large\textbf{CNR-IEIIT}}\\
Istituto di Elettronica e di Ingegneria dell'Informazione e delle Telecomunicazioni\\
Consiglio Nazionale delle Ricerche, Milan, Italy\\[1.5em]
% ----------------------------------------------------------------------------

\begin{center}
\textbf{\Large Letter of Intent (LoI)}
\end{center}

\bigskip
\begin{flushright}
Mr.\ Robin Wydaeghe and Prof.\ Wout Joseph\\
Ghent University -- imec\\[1em]
Milan, 22 June 2026
\end{flushright}

\bigskip
\textbf{Subject: Letter of Intent (LoI) regarding the fast numerical-dosimetry tool (AEGIS),
by Robin Wydaeghe}

\bigskip
Dear Mr.\ Wydaeghe, dear Prof.\ Joseph,

As Director of Research at CNR-IEIIT in Milan, working on numerical and stochastic dosimetry
of human exposure to electromagnetic fields, I am very aware of how much the computational
cost of full-wave (FDTD) simulation limits what we can study in practice. A method that is
very substantially faster than FDTD with comparable accuracy, and that is differentiable,
would be of real value to our work.

In our high-frequency work the mesh is the constant bottleneck, \textbf{it is our `nightmare'},
and the time my researchers spend waiting for FDTD simulations to finish amounts to
\textbf{over 1.5 FTE}. As I put it during our conversation, \textbf{if I find a real improvement
over what I have now, I would be able to pay for it}.

Should this method become available as a software tool, CNR-IEIIT would therefore be
interested, in principle, in evaluating it and, if it fits our needs, in licensing it for use
in our research group.

This is a letter of intent. It is entirely non-binding, creates no obligation on either side,
and any future use would naturally depend on availability, terms, and our needs at the time.
The practical needs and constraints behind this interest are set out in more detail in the
appended record of our interview of 22 June 2026 (Appendix~A), which forms part of this letter.

On this basis I am glad to support Robin Wydaeghe's IOF proposal, and I look forward to seeing
the tool take shape.

\bigskip
With best regards,

\vspace{2.2cm}

\rule{6cm}{0.4pt}\\
Dr.\ Marta Parazzini\\
Director of Research, CNR-IEIIT\\
Consiglio Nazionale delle Ricerche, Milan, Italy\\
marta.parazzini@cnr.it \hfill Date: 22 June 2026

% ============================ APPENDIX ======================================
\newpage

\begin{center}
\textbf{\large Appendix A --- Record of interview (22 June 2026)}\\[0.3em]
Appendix to the Letter of Intent above, in support of the IOF valorisation project (AEGIS)
\end{center}

\bigskip
On 22 June 2026, I had a video conference call with Mr.\ Robin Wydaeghe. I was asked
questions in a non-leading, objective way. During this call, I am either quoted or faithfully
paraphrased in the following statements:

\medskip
\begin{itemize}[leftmargin=1.4em,itemsep=0.5em]

  \item There are 4 to 5 researchers working with FDTD in my research group. Two of them work
  in the high-frequency regime ($>6$\,GHz), three of them in the low-frequency regime for
  biomedical applications ($<6$\,GHz).

  \item With the advent of 6G comprising the new Frequency Band 3, I intend to hire more
  researchers working in the high-frequency regime in the coming years.

  \item Each of the employees in this category spends around 30--50\% of their time waiting
  for FDTD simulations to finish, which translates to \textbf{1.5 FTE} in total (rising in the future).
  The rest of their time is spent on data analysis, paper writing, reading literature, and
  similar tasks.

  \item Configuring FDTD simulations takes around 10\% of a senior researcher's time and
  around 30\% of a junior researcher's time.

  \item The bottleneck in their research is always the mesh; \textbf{it is our `nightmare'}.

  \item Researchers perform high-frequency simulations in both the far-field (incident fields
  in an environment) and the near-field (incident from a device).

  \item Three things would help convince me to use a fast numerical dosimetry method, or to
  recommend it to others:
  \begin{enumerate}[leftmargin=1.6em,topsep=0.3em,itemsep=0.2em]
    \item Clear tutorials and an easy-to-use platform are extremely useful. This enables us to
    validate it ourselves while learning it. I rate the importance of this 5/5.
    \item A peer-reviewed paper, in the sense that it validates the method. I rate the
    importance of this 4/5.
    \item If I am performing standards-related work, it is useful if it is suggested as a
    viable method in the standards. I rate the importance of this 3/5.
  \end{enumerate}
  Point 1 is likely enough to convince me, with the others strengthening my decision.

  \item On the commercial value of the method: \textbf{``If I really find a real improvement
  with respect to what I have now, I would be able to pay you.''}

  \item Currently, we own 4 licences necessary to run high-frequency exposure simulations,
  each at 3\,kEUR per year, and various other ones (such as those for phantoms) similarly
  priced.

  \item Zurich MedTech provides some licences for free as part of their University Research
  Agreement, but only when we use them for research purposes. We also pay for discounted CST
  licences.

\end{itemize}

\vspace{1.5em}
\rule{6cm}{0.4pt}\\
Dr.\ Marta Parazzini, Director of Research, CNR-IEIIT \hfill Date: 22 June 2026

\end{document}
